% ============================================================
% f2014_outline_arts_q19_prism.tex
% 2014 全国大纲卷（文科）第 19 题：三棱柱 ABC-A1B1C1
% A1 在平面 ABC 内的射影 D 在 AC 上，∠ACB=90°。
% 来源 PDF：普通高考/2014/2014大纲文(广西).pdf
% 原页：第 1 页（A3 横排）第 3 栏，印面 871，第 19 题题干下方
% ============================================================
\begin{tikzpicture}[>=Stealth, line width=0.7pt, font=\small,
  line cap=round, line join=round]
  \coordinate (A) at (0,0);
  \coordinate (C) at (2.55,0.85);
  \coordinate (B) at (4.35,0.65);
  \coordinate (D) at (0.95,0.3167);
  \coordinate (A1) at (0.72,3.15);
  \coordinate (C1) at (3.27,4.00);
  \coordinate (B1) at (5.07,3.80);

  % 可见棱
  \draw (A) -- (B) -- (B1) -- (C1) -- (A1) -- cycle;
  \draw (A) -- (A1);
  \draw (A1) -- (B1);
  \draw (A1) -- (B);

  % 不可见棱与辅助线
  \draw[dashed] (A) -- (C) -- (B);
  \draw[dashed] (C) -- (C1);
  \draw[dashed] (A1) -- (D);
  \draw[dashed] (A1) -- (C);

  \fill (A) circle (1.1pt) node[below left] {$A$};
  \fill (B) circle (1.1pt) node[right] {$B$};
  \fill (C) circle (1.1pt) node[above left] {$C$};
  \fill (D) circle (1.1pt) node[below right] {$D$};
  \fill (A1) circle (1.1pt) node[left] {$A_1$};
  \fill (C1) circle (1.1pt) node[above] {$C_1$};
  \fill (B1) circle (1.1pt) node[above right] {$B_1$};
\end{tikzpicture}
